\ifajar
  \section{Latihan}
\else
  \section{Bab 12: Type Safety, Kontrak API, dan Desain Layanan}
\fi

Seluruh file yang dipakai pada latihan ini tersedia di
\berkas{sources/chapter12/}, dihitung relatif terhadap akar repositori.
Latihan ini membutuhkan Python 3 beserta \texttt{jsonschema}, dan
\texttt{tsc} yang sudah terpasang di Bab 6. Tidak ada Docker dan tidak ada
pemasangan npm baru.

Seluruh permintaan diarahkan ke aplikasi dan \textit{gateway} yang berjalan
di komputer sendiri, dan sifatnya mengamati, bukan membebani. Status dan
klasifikasi perubahan akan sama pada tiap mahasiswa, sedangkan lama
pengukuran dan nilai \texttt{Retry-After} akan berbeda sedikit.

Penyiapannya mengaktifkan venv bersama di \berkas{sources/.venv}, yang memuat
\texttt{jsonschema}, lalu menyalakan dua layanan, seperti pada Listing
\ref{lst:c12-siapkan}.

\begin{lstlisting}[language=bash, caption={Menyalakan aplikasi dan gateway, tersedia di \berkas{sources/chapter12/}}, label={lst:c12-siapkan}]
cd sources/chapter12
source ../.venv/bin/activate
python3 aplikasi.py &
python3 gateway.py &

# Keluarannya:
# Melayani 48 buku di http://127.0.0.1:8090/v1/buku  (Ctrl+C untuk berhenti)
# Gateway di http://127.0.0.1:8091/  (Ctrl+C untuk berhenti)
# Ember 20 token per klien, diisi 5 token per detik
\end{lstlisting}

\ifajar
  \subsection*{Latihan 1: Meng-generate Tipe dari Kontrak}
\else
  \subsection*{Prosedur 1: Meng-generate Tipe dari Kontrak}
\fi

Latihan ini membuktikan bahwa tipe klien dapat diturunkan dari kontrak, dan
bahwa perubahan kontrak tertangkap compiler sebelum program jalan.

\begin{enumerate}
  \item Jalankan generatornya untuk kedua kontrak, seperti pada Listing
        \ref{lst:c12-lat1-jalan}, lalu catat jumlah skema dan jumlah baris
        hasilnya.
  \item Periksa seluruh tipe dengan \texttt{tsc} memakai
        \berkas{tsconfig.json}, lalu catat hasilnya.
  \item Periksa \berkas{pakai-tipe-salah.ts} sendiri, lalu catat kode
        kesalahan beserta nomor barisnya.
  \item Tambahkan satu field opsional pada skema \texttt{Buku} di
        \berkas{kontrak.json}, jalankan ulang generatornya, lalu catat
        perubahan pada file tipenya.
\ifajar
  \item Isi \tabisi{tab:c12-lat1-hasil}. Tabel \ref{tab:c12-lat1-contoh}
        memperlihatkan contoh yang sudah terisi.
\else
  \item Bandingkan hasil pengukuran sendiri dengan Tabel
        \ref{tab:c12-lat1-contoh}.
\fi
  \item Jawab dua pertanyaan berikut. Pertama, mengapa field yang tidak ada
        di \texttt{required} di-\textit{generate} dengan tanda tanya? Kedua,
        kesalahan apa yang tetap lolos dari \texttt{tsc} tetapi tertangkap
        validasi di server?
\end{enumerate}

\begin{lstlisting}[language=bash, caption={Meng-generate tipe lalu memeriksanya, tersedia di \berkas{sources/chapter12/buat\_tipe.py}}, label={lst:c12-lat1-jalan}]
python3 buat_tipe.py kontrak.json tipe-katalog.ts
python3 buat_tipe.py kontrak-v2.json tipe-katalog-v2.ts
../chapter06/node_modules/.bin/tsc --noEmit -p tsconfig.json

# Keluarannya:
# 5 skema dan 2 jalur di-generate menjadi tipe-katalog.ts
# 5 skema dan 2 jalur di-generate menjadi tipe-katalog-v2.ts
\end{lstlisting}

\begin{table}[!htb]
  \centering
  \caption{Contoh hasil Latihan 1 yang sudah terisi}
  \label{tab:c12-lat1-contoh}
  \begin{tabularx}{\textwidth}{Xll}
    \hline
    \textbf{Yang diperiksa} & \textbf{Hasil} & \textbf{Keterangan} \\
    \hline
    Jumlah skema dan jalur & 5 dan 2 & dari \berkas{kontrak.json} \\
    \hline
    Baris tipe hasil generate & 50 & versi 2 menjadi 51 \\
    \hline
    \texttt{tsc} pada seluruh tipe & lolos & tanpa kesalahan \\
    \hline
    \texttt{tsc} pada file salah & TS2339 & baris 15, field \texttt{harga} \\
    \hline
  \end{tabularx}
\end{table}

\ifajar
\begin{table}[!htb]
  \centering
  \caption{Lembar isian hasil pengamatan Latihan 1}
  \label{tab:c12-lat1-hasil}
  \begin{tabularx}{\textwidth}{Xll}
    \hline
    \textbf{Yang diperiksa} & \textbf{Hasil} & \textbf{Keterangan} \\
    \hline
    Jumlah skema dan jalur & \dots & \dots \\
    \hline
    Baris tipe hasil generate & \dots & \dots \\
    \hline
    \texttt{tsc} pada seluruh tipe & \dots & \dots \\
    \hline
    \texttt{tsc} pada file salah & \dots & \dots \\
    \hline
  \end{tabularx}
\end{table}
\fi

\ifajar
  \subsection*{Latihan 2: Validasi di Batas dan Idempotensi}
\else
  \subsection*{Prosedur 2: Validasi di Batas dan Idempotensi}
\fi

Latihan ini membuktikan bahwa satu pemeriksaan di pintu masuk cukup untuk
menolak lima bentuk isi yang salah, dan bahwa pengiriman ulang tidak membuat
pesanan kedua.

\begin{enumerate}
  \item Nyalakan ulang aplikasinya, agar nomor pesanannya dimulai dari awal.
  \item Jalankan \berkas{uji-kontrak.py}, seperti pada Listing
        \ref{lst:c12-lat2-jalan}, lalu catat sebelas barisnya.
  \item Kirim satu permintaan dengan kunci idempotensi pilihan sendiri dua
        kali, lalu catat kedua statusnya dan nomor pesanannya.
  \item Kirim permintaan ketiga dengan kunci yang sama tetapi
        \texttt{id\_buku} berbeda, lalu catat statusnya.
\ifajar
  \item Isi \tabisi{tab:c12-lat2-hasil}. Tabel \ref{tab:c12-lat2-contoh}
        memperlihatkan contoh yang sudah terisi.
\else
  \item Bandingkan hasil pengukuran sendiri dengan Tabel
        \ref{tab:c12-lat2-contoh}.
\fi
  \item Jawab tiga pertanyaan berikut. Pertama, mengapa pengiriman kedua
        dijawab 200, bukan 201? Kedua, mengapa field asing ditolak, bukan
        diabaikan? Ketiga, apa yang terjadi bila kunci idempotensinya dibuat
        dari isi permintaan, bukan dibuat klien?
\end{enumerate}

\begin{lstlisting}[language=bash, caption={Menguji validasi dan idempotensi, tersedia di \berkas{sources/chapter12/uji-kontrak.py}}, label={lst:c12-lat2-jalan}]
python3 uji-kontrak.py

# Keluarannya:
#   Yang diuji                   Status  Hasil
#   GET /v1/buku                    200  sesuai skema
#   GET /v1/buku?halaman=2          200  sesuai skema
#   GET /v1/buku?per_halaman=99     422  sesuai skema
#   sesuai kontrak                  201  sesuai, diterima
#   field wajib hilang              422  sesuai, isi tidak sesuai
#   tipe salah                      422  sesuai, isi tidak sesuai
#   melebihi batas                  422  sesuai, isi tidak sesuai
#   field asing                     422  sesuai, isi tidak sesuai
#   kunci baru                      201  pesanan psn-0002 dibuat
#   kunci sama, isi sama            200  pesanan psn-0002 dikembalikan, sama
#   kunci sama, isi beda            409  kunci idempotensi dipakai ulang
\end{lstlisting}

\begin{table}[!htb]
  \centering
  \caption{Contoh hasil Latihan 2 yang sudah terisi}
  \label{tab:c12-lat2-contoh}
  \begin{tabularx}{\textwidth}{Xll}
    \hline
    \textbf{Yang dikirim} & \textbf{Status} & \textbf{Hasil} \\
    \hline
    isi sesuai kontrak & 201 & diterima \\
    \hline
    field wajib hilang & 422 & isi tidak sesuai \\
    \hline
    field asing \texttt{catatan} & 422 & isi tidak sesuai \\
    \hline
    kunci sama, isi sama & 200 & pesanan yang sama dikembalikan \\
    \hline
    kunci sama, isi beda & 409 & ditolak \\
    \hline
  \end{tabularx}
\end{table}

\ifajar
\begin{table}[!htb]
  \centering
  \caption{Lembar isian hasil pengamatan Latihan 2}
  \label{tab:c12-lat2-hasil}
  \begin{tabularx}{\textwidth}{Xll}
    \hline
    \textbf{Yang dikirim} & \textbf{Status} & \textbf{Hasil} \\
    \hline
    isi sesuai kontrak & \dots & \dots \\
    \hline
    field wajib hilang & \dots & \dots \\
    \hline
    field asing \texttt{catatan} & \dots & \dots \\
    \hline
    kunci sama, isi sama & \dots & \dots \\
    \hline
    kunci sama, isi beda & \dots & \dots \\
    \hline
  \end{tabularx}
\end{table}
\fi

\ifajar
  \subsection*{Latihan 3: Memisahkan Perubahan Aman dari yang Merusak}
\else
  \subsection*{Prosedur 3: Memisahkan Perubahan Aman dari yang Merusak}
\fi

Latihan ini membuktikan bahwa jenis perubahan kontrak dapat diputuskan mesin,
dan bahwa klien lama dapat rusak walaupun server menjawab 200.

\begin{enumerate}
  \item Jalankan \berkas{uji-kompatibilitas.py}, seperti pada Listing
        \ref{lst:c12-lat3-jalan}, lalu catat lima perubahan beserta
        akibatnya.
  \item Catat keempat baris pengujian pada server, terutama statusnya.
  \item Tambahkan satu field opsional pada skema \texttt{Buku} di
        \berkas{kontrak-v2.json}, jalankan ulang, lalu catat perubahannya di
        tabel perbandingan.
  \item Ganti nama satu field wajib di \berkas{kontrak-v2.json}, jalankan
        ulang, lalu catat berapa baris "merusak klien lama" yang muncul.
\ifajar
  \item Isi \tabisi{tab:c12-lat3-hasil}. Tabel \ref{tab:c12-lat3-contoh}
        memperlihatkan contoh yang sudah terisi.
\else
  \item Bandingkan hasil pengukuran sendiri dengan Tabel
        \ref{tab:c12-lat3-contoh}.
\fi
  \item Jawab tiga pertanyaan berikut. Pertama, mengapa mengganti nama satu
        field menghasilkan dua baris perubahan? Kedua, mengapa menambah nilai
        enum merusak klien lama? Ketiga, apa yang harus disiapkan sebelum
        versi 1 dimatikan?
\end{enumerate}

\begin{lstlisting}[language=bash, caption={Membandingkan kontrak lalu mengujinya, tersedia di \berkas{sources/chapter12/uji-kompatibilitas.py}}, label={lst:c12-lat3-jalan}]
python3 uji-kompatibilitas.py

# Keluarannya:
#   Perubahan kontrak v1 ke v2                     Akibat
#   Buku.harga wajib dihapus                       merusak klien lama
#   Buku.harga_rupiah wajib ditambah               merusak klien lama
#   Buku.penerbit opsional ditambah                aman
#   Pesanan.status enum menerima dibatalkan        merusak klien lama
#   PesananBaru.jumlah batas naik ke 20            aman bagi klien lama
#
#   Yang diuji pada server       Status  Hasil
#   klien v1 membaca /v1/buku       200  jalan
#   klien v1 membaca /v2/buku       200  gagal: 'harga' is a required property
#   klien v2 memesan 15 ke /v2      201  diterima
#   klien v2 memesan 15 ke /v1      422  isi tidak sesuai
\end{lstlisting}

\begin{table}[!htb]
  \centering
  \caption{Contoh hasil Latihan 3 yang sudah terisi}
  \label{tab:c12-lat3-contoh}
  \begin{tabularx}{\textwidth}{Xll}
    \hline
    \textbf{Yang diamati} & \textbf{Hasil} & \textbf{Keterangan} \\
    \hline
    Perubahan terdeteksi & 5 & 2 aman, 3 merusak \\
    \hline
    Klien v1 ke \texttt{/v2/buku} & 200 & gagal dipakai \\
    \hline
    Klien v2 ke \texttt{/v1/pesanan} & 422 & batas lama 10 \\
    \hline
    Ganti nama satu field wajib & 2 baris & dihapus dan ditambah \\
    \hline
  \end{tabularx}
\end{table}

\ifajar
\begin{table}[!htb]
  \centering
  \caption{Lembar isian hasil pengamatan Latihan 3}
  \label{tab:c12-lat3-hasil}
  \begin{tabularx}{\textwidth}{Xll}
    \hline
    \textbf{Yang diamati} & \textbf{Hasil} & \textbf{Keterangan} \\
    \hline
    Perubahan terdeteksi & \dots & \dots \\
    \hline
    Klien v1 ke \texttt{/v2/buku} & \dots & \dots \\
    \hline
    Klien v2 ke \texttt{/v1/pesanan} & \dots & \dots \\
    \hline
    Ganti nama satu field wajib & \dots & \dots \\
    \hline
  \end{tabularx}
\end{table}
\fi

\ifajar
  \subsection*{Latihan 4: Batas Laju di Gateway}
\else
  \subsection*{Prosedur 4: Batas Laju di Gateway}
\fi

Latihan ini mengukur perilaku \textit{token bucket}: lonjakan pendek yang
masih dilayani, penolakan sesudah embernya habis, dan pengisian ulangnya.

\begin{enumerate}
  \item Pastikan aplikasi dan \textit{gateway} menyala, lalu jalankan
        \berkas{ukur-gateway.py} seperti pada Listing
        \ref{lst:c12-lat4-jalan}.
  \item Catat jumlah 200 dan 429 pada ketiga tahapnya, beserta
        \texttt{Retry-After} yang diterima.
  \item Hitung sendiri jumlah yang seharusnya dilayani pada tahap kedua dari
        lama menunggu dan laju pengisiannya, lalu bandingkan dengan hasil
        pengukuran.
  \item Ubah \texttt{KAPASITAS\_EMBER} menjadi 5 di \berkas{gateway.py},
        nyalakan ulang \textit{gateway}-nya, lalu ukur sekali lagi.
\ifajar
  \item Isi \tabisi{tab:c12-lat4-hasil}. Tabel \ref{tab:c12-lat4-contoh}
        memperlihatkan contoh yang sudah terisi.
\else
  \item Bandingkan hasil pengukuran sendiri dengan Tabel
        \ref{tab:c12-lat4-contoh}.
\fi
  \item Jawab tiga pertanyaan berikut. Pertama, mengapa klien lain tetap
        dilayani 20 permintaan? Kedua, apa akibatnya bila embernya dibuat
        satu untuk seluruh sistem? Ketiga, mengapa penolakannya memakai 429
        beserta \texttt{Retry-After}, bukan 503 tanpa keterangan?
\end{enumerate}

\begin{lstlisting}[language=bash, caption={Mengukur batas laju gateway, tersedia di \berkas{sources/chapter12/ukur-gateway.py}}, label={lst:c12-lat4-jalan}]
python3 ukur-gateway.py

# Keluarannya:
#   Tahap                            200   429  Retry-After  Lama detik
#   30 permintaan sekaligus           20    10         0.19        0.02
#   setelah menunggu 2 detik          10    20         0.16        0.02
#   klien lain, ember sendiri         20    10         0.19        0.02
\end{lstlisting}

\begin{table}[!htb]
  \centering
  \caption{Contoh hasil Latihan 4 yang sudah terisi}
  \label{tab:c12-lat4-contoh}
  \begin{tabularx}{\textwidth}{Xlll}
    \hline
    \textbf{Tahap} & \textbf{200} & \textbf{429} & \textbf{Hitungan sendiri} \\
    \hline
    30 sekaligus & 20 & 10 & kapasitas ember 20 \\
    \hline
    setelah 2 detik & 10 & 20 & $2 \times 5$ token \\
    \hline
    klien lain & 20 & 10 & ember terpisah \\
    \hline
  \end{tabularx}
\end{table}

\ifajar
\begin{table}[!htb]
  \centering
  \caption{Lembar isian hasil pengukuran Latihan 4}
  \label{tab:c12-lat4-hasil}
  \begin{tabularx}{\textwidth}{Xlll}
    \hline
    \textbf{Tahap} & \textbf{200} & \textbf{429} & \textbf{Hitungan sendiri} \\
    \hline
    30 sekaligus & \dots & \dots & \dots \\
    \hline
    setelah 2 detik & \dots & \dots & \dots \\
    \hline
    klien lain & \dots & \dots & \dots \\
    \hline
  \end{tabularx}
\end{table}
\fi
